\chapter{Determinação estrutural: RMN de \texorpdfstring{$^{13}$C}{13C}, EM e IV combinados}\label{ch:b2:structure-determination}

Um laboratório de perfumaria recebe um frasco de um líquido incolor com cheiro de jasmim. Ninguém
escreveu um rótulo. Um \omterm{def:b2:mass-spec-atomic:spectrum}{espectro de massas} dá sua fórmula molecular, um espectro infravermelho seus grupos
funcionais, um espectro de carbono 13 seu esqueleto e um espectro de prótons a maneira como as peças se
unem; em uma hora o líquido tem um nome. Este capítulo acrescenta a \omterm{def:b2:structure-determination:c13}{RMN de carbono 13} à RMN de prótons e
à espectroscopia infravermelha do volume do 1.º ano e à \omterm{def:b2:mass-spec-atomic:spectrum}{espectrometria de massas} do
\cref{ch:b2:mass-spec-atomic}, e transforma as quatro em um único método.

\begin{recall}
O volume do 1.º ano: deslocamento químico e blindagem, prótons equivalentes, integração, acoplamento
spin--spin e a regra do $n + 1$, números de onda dos grupos funcionais no infravermelho e a região da
impressão digital, o grau de insaturação, a resolução de uma estrutura a partir de três espectros.
\Cref{ch:b2:mass-spec-atomic}: o \omterm{def:b2:mass-spec-atomic:molecular-ion}{íon molecular}, os \omterm{def:b2:mass-spec-atomic:isotope-pattern}{perfis isotópicos}, a fragmentação.
\end{recall}

\section{A RMN de carbono 13}

O carbono 12, o isótopo principal, não tem spin nuclear e não dá sinal de RMN. O carbono 13 tem spin
$\tfrac12$, como o próton, mas representa apenas 1.1~\% do carbono natural e seu sinal é mais fraco,
núcleo por núcleo, que o de um próton; um espectro de carbono exige portanto mais amostra ou mais
varreduras.
% ledger: iso:C13

\begin{definition}[RMN de carbono 13]\label{def:b2:structure-determination:c13}
A \emph{RMN de carbono 13}\index{RMN de carbono 13} registra as ressonâncias dos núcleos \ce{^{13}C} de uma amostra.
Ela é em geral realizada com \emph{desacoplamento de banda larga}\index{desacoplamento de banda larga}: os prótons são
irradiados em toda a sua faixa de frequências durante a medida, o que remove todo acoplamento
carbono--próton, de modo que cada tipo de carbono dá uma única raia. \emph{Carbonos
equivalentes}\index{carbonos equivalentes}, permutados por uma simetria da molécula (ou por uma rotação rápida), dão a
mesma raia.
\end{definition}

\begin{proposition}[Nenhum acoplamento carbono--carbono]\label{prop:b2:structure-determination:no-cc-coupling}
Os acoplamentos entre átomos de carbono vizinhos não aparecem em um espectro de \ce{^{13}C} em
abundância natural.
\end{proposition}

\begin{proof}
Um acoplamento entre dois carbonos só é visto nas moléculas em que ambos são \ce{^{13}C}. Para duas
posições vizinhas dadas, a probabilidade é $a_{13}^2 = 0.011^2 \approx 1.2 \times 10^{-4}$: cerca de uma
molécula em oito mil. O sinal de cada carbono vem quase inteiramente de moléculas em que seus vizinhos são
\ce{^{12}C}, invisíveis para o espectrômetro; os satélites acoplados são cem vezes mais fracos que o ruído
de um espectro comum.
\end{proof}

\begin{proposition}[Número de raias]\label{prop:b2:structure-determination:signals}
Um espectro de \ce{^{13}C} desacoplado mostra uma raia por conjunto de \omterm{def:b2:structure-determination:c13}{carbonos equivalentes}, a menos que
dois conjuntos tenham por acaso o mesmo deslocamento.
\end{proposition}

\begin{proof}[Raciocínio]
\omterm{def:b2:structure-determination:c13}{Carbonos equivalentes} estão em ambientes idênticos, portanto têm o mesmo deslocamento; carbonos não
equivalentes em geral diferem. Com os acoplamentos com os prótons removidos e os acoplamentos entre
carbonos ausentes (\cref{prop:b2:structure-determination:no-cc-coupling}), nada desdobra uma raia. A
sobreposição acidental ocorre quando dois carbonos diferentes têm deslocamentos mais próximos que a
resolução, o que é comum para os carbonos de um anel benzênico perto de \qty{128}{ppm}.
\end{proof}

Os deslocamentos se espalham por cerca de \qty{220}{ppm}, vinte vezes a faixa dos prótons, de modo que
as sobreposições são mais raras que nos espectros de prótons, e o deslocamento sozinho indica o tipo de
carbono. O gráfico abaixo mostra os deslocamentos medidos para os compostos deste capítulo.

\begin{omfigure}
\begin{tikzpicture}
\begin{axis}[width=0.86\linewidth, height=4.6cm, xmin=0, xmax=220, ymin=-0.6, ymax=4.6, x dir=reverse,
  xlabel={$\delta$ (ppm)}, label style={font=\small}, tick label style={font=\scriptsize},
  ytick={0,1,2,3,4}, yticklabels={C alquila, C--O, C aromático, C=O de éster, C=O de cetona},
  xtick={0,20,40,60,80,100,120,140,160,180,200,220}, grid=major, grid style={black!10}]
\addplot[only marks, mark=|, mark size=5pt, thick, omDef] table[x=shift, y=cls] {figdata/out/bachelor-2-nmr-spectra-d.dat};
\end{axis}
\end{tikzpicture}

{\small Deslocamentos de \ce{^{13}C} medidos da butan-2-ona, da pentan-2-ona, do benzoato de etila e do
etanoato de benzila, classificados por tipo de carbono (CDCl$_3$; dados do PubChem). As carbonilas de
cetona ficam perto de 209, as de éster perto de 167--171, os carbonos \omterm{def:b2:huckel:aromatic}{aromáticos} entre 128 e 137, os
carbonos ligados a um oxigênio de éster perto de 61--66, e os carbonos alquila abaixo de 46.}
% figdata: bachelor-2/nmr-spectra.py part d
% ledger: c13:butanone, c13:pentan-2-one, c13:ethylbenzoate, c13:benzylacetate
\end{omfigure}

\begin{proposition}[Alturas não são contagens]\label{prop:b2:structure-determination:intensities}
Em um espectro de \ce{^{13}C} desacoplado de rotina, as alturas das raias não são proporcionais aos
números de carbonos que elas representam; os carbonos sem hidrogênio dão raias fracas.
\end{proposition}

\begin{proof}[Raciocínio]
Os espectros de rotina repetem a medida mais depressa do que os carbonos mais lentos relaxam de volta ao
equilíbrio, e os carbonos sem prótons ligados relaxam lentamente, de modo que seu sinal fica parcialmente
saturado. O desacoplamento também intensifica os sinais dos carbonos que levam prótons, por uma
quantidade que depende de seu ambiente. Na butan-2-ona, a raia da carbonila é a mais fraca das quatro,
embora represente um carbono como as outras. A integração, tão útil para os prótons, não é portanto
usada nos espectros de carbono de rotina.
\end{proof}

\section{DEPT}

\begin{definition}[DEPT]\label{def:b2:structure-determination:dept}
O \emph{DEPT}\index{DEPT} (\textit{distortionless enhancement by polarisation transfer}) é um conjunto de
experimentos de carbono 13 que dão aos sinais dos carbonos com prótons ligados um sinal ou uma ausência
conforme seu número de hidrogênios. Um \emph{carbono quaternário}\index{carbono quaternario@carbono quaternário} não leva
hidrogênio (ligado a quatro outros átomos, ou um carbono carbonílico ou \omterm{def:b2:huckel:aromatic}{aromático} sem H); ele não dá
sinal em DEPT.
\end{definition}

\begin{proposition}[Ler um DEPT]\label{prop:b2:structure-determination:dept}
Em DEPT-135, os carbonos \ce{CH} e \ce{CH3} apontam para cima, os carbonos \ce{CH2} para baixo e os
\omterm{def:b2:structure-determination:dept}{carbonos quaternários} estão ausentes. Em DEPT-90, só aparecem os carbonos \ce{CH}.
\end{proposition}

\admitted

O experimento transfere a magnetização dos prótons para o carbono ao qual estão ligados, por meio de uma
sequência de pulsos de radiofrequência cujo último ângulo seleciona a resposta; como os pulsos o fazem é
explicado no volume do 3.º ano. Comparar o espectro desacoplado com DEPT-135 e DEPT-90 classifica cada
raia em C, CH, \ce{CH2} ou \ce{CH3}.

\begin{omfigure}
\pgfplotsset{dept/.style={width=\linewidth, height=2.9cm, ycomb, x dir=reverse, ymin=-1, ymax=1.15,
  ytick=\empty, axis y line=none, axis x line=middle, x axis line style={-}, tick label style={font=\tiny},
  title style={font=\scriptsize, at={(0.0,0.95)}, anchor=north west}, every axis plot/.append style={thick}}}
\begin{minipage}[t]{0.49\linewidth}\centering
{\small butan-2-ona}\par
\begin{tikzpicture}
\begin{axis}[dept, xmin=0, xmax=220, xtick={0,50,100,150,200}, title={desacoplado}]
\addplot[omDef] table[x=shift, y=dec] {figdata/out/bachelor-2-nmr-spectra-a.dat};
\end{axis}
\end{tikzpicture}\par
\begin{tikzpicture}
\begin{axis}[dept, xmin=0, xmax=220, xtick={0,50,100,150,200}, title={DEPT-135}]
\addplot[omThm] table[x=shift, y=d135] {figdata/out/bachelor-2-nmr-spectra-a.dat};
\end{axis}
\end{tikzpicture}\par
\begin{tikzpicture}
\begin{axis}[dept, xmin=0, xmax=220, xtick={0,50,100,150,200}, title={DEPT-90}]
\addplot[omMeth] table[x=shift, y=d90] {figdata/out/bachelor-2-nmr-spectra-a.dat};
\end{axis}
\end{tikzpicture}
\end{minipage}\hfill
\begin{minipage}[t]{0.49\linewidth}\centering
{\small benzoato de etila}\par
\begin{tikzpicture}
\begin{axis}[dept, xmin=0, xmax=220, xtick={0,50,100,150,200}, title={desacoplado}]
\addplot[omDef] table[x=shift, y=dec] {figdata/out/bachelor-2-nmr-spectra-b.dat};
\end{axis}
\end{tikzpicture}\par
\begin{tikzpicture}
\begin{axis}[dept, xmin=0, xmax=220, xtick={0,50,100,150,200}, title={DEPT-135}]
\addplot[omThm] table[x=shift, y=d135] {figdata/out/bachelor-2-nmr-spectra-b.dat};
\end{axis}
\end{tikzpicture}\par
\begin{tikzpicture}
\begin{axis}[dept, xmin=0, xmax=220, xtick={0,50,100,150,200}, title={DEPT-90}]
\addplot[omMeth] table[x=shift, y=d90] {figdata/out/bachelor-2-nmr-spectra-b.dat};
\end{axis}
\end{tikzpicture}
\end{minipage}

{\small Espectros de \ce{^{13}C} desacoplados (deslocamentos e alturas medidos, dados do PubChem) com as
respostas \omterm{def:b2:structure-determination:dept}{DEPT} (sinais segundo a \cref{prop:b2:structure-determination:dept}, alturas esquemáticas).
Butan-2-ona: C=O 209.3 (ausente em \omterm{def:b2:structure-determination:dept}{DEPT}), \ce{CH2} 36.9 (para baixo), \ce{CH3} 29.4 e 7.9 (para cima).
Benzoato de etila: C=O 166.5 e o carbono do anel que leva o éster, 130.6, desaparecem; três CH \omterm{def:b2:huckel:aromatic}{aromáticos}
(132.8, 129.6, 128.3) permanecem em DEPT-90; \ce{OCH2} 60.9 aponta para baixo.}
% figdata: bachelor-2/nmr-spectra.py parts a, b
\end{omfigure}

\section{Combinar as técnicas}

Cada técnica responde à sua própria pergunta, e as respostas são montadas em uma ordem fixa: a fórmula
primeiro, porque ela limita todo o resto; depois os grupos funcionais e o esqueleto; depois as
conexões.

\begin{method}[Resolver uma substância desconhecida]\label{met:b2:structure-determination:solve}
\begin{enumerate}
  \item Fórmula: \omterm{def:b2:mass-spec-atomic:molecular-ion}{íon molecular}, $M+1$ para o número de carbonos, \omterm{def:b2:mass-spec-atomic:isotope-pattern}{perfis isotópicos}, \omterm{def:b2:mass-spec-atomic:nitrogen-rule}{regra do nitrogênio};
        uma massa exata se disponível. Grau de insaturação.
  \item Infravermelho: \ce{C=O} (e de que tipo), \ce{O-H}, \ce{N-H}, \ce{C#N}, \ce{C-H} \omterm{def:b2:huckel:aromatic}{aromático} ou de
        alceno acima de \qty{3000}{cm^{-1}}.
  \item Carbono 13 e \omterm{def:b2:structure-determination:dept}{DEPT}: número de tipos de carbono (simetria), suas classes pelos deslocamentos, e
        C, CH, \ce{CH2}, \ce{CH3} pelo \omterm{def:b2:structure-determination:dept}{DEPT}. Compare com a fórmula.
  \item RMN de prótons: integrações, deslocamentos, multiplicidades; vizinhos pela regra do $n + 1$.
  \item Monte os fragmentos; escreva todas as estruturas candidatas.
  \item Confronte cada candidata com cada dado, incluindo os fragmentos do \omterm{def:b2:mass-spec-atomic:spectrum}{espectro de massas}; fique com
        aquela que explica todos.
\end{enumerate}
\end{method}

\begin{omfigure}
\begin{tikzpicture}[x=1cm, y=1cm, st/.style={draw=black!70, rounded corners=2pt, fill=white,
  font=\small, align=center, minimum height=0.9cm, inner sep=3pt, text width=2.3cm}]
\node[st, fill=omDef!10] (ms) at (0,0) {EM\\fórmula, $M+1$};
\node[st] (dbe) at (2.9,0) {grau de\\insaturação};
\node[st, fill=omDef!10] (ir) at (5.8,0) {IV\\grupos funcionais};
\node[st, fill=omDef!10] (c13) at (8.7,0) {\ce{^{13}C}, DEPT\\esqueleto};
\node[st, fill=omDef!10] (h1) at (8.7,-1.8) {\ce{^1H}\\conexões};
\node[st] (cand) at (5.8,-1.8) {estruturas\\candidatas};
\node[st, fill=omThm!10] (chk) at (2.9,-1.8) {confrontar com\\cada dado};
\node[st] (ans) at (0,-1.8) {estrutura};
\draw[->, thick] (ms) -- (dbe); \draw[->, thick] (dbe) -- (ir); \draw[->, thick] (ir) -- (c13);
\draw[->, thick] (c13) -- (h1); \draw[->, thick] (h1) -- (cand); \draw[->, thick] (cand) -- (chk);
\draw[->, thick] (chk) -- (ans);
\draw[->, thick, densely dashed] (chk.south) -- ++(0,-0.5) -| node[pos=0.25, below, font=\footnotesize]
  {um dado não explicado: de volta às candidatas} (cand.south);
\end{tikzpicture}
\par\vspace{4pt}

{\small A ordem de trabalho para uma substância desconhecida (\cref{met:b2:structure-determination:solve}).
A última etapa é a mais frequentemente omitida: uma estrutura que deixa um pico sem explicação ainda não
é a resposta.}
\end{omfigure}

\section{Casos resolvidos}

\begin{example}[Uma cetona, C$_5$H$_{10}$O]\label{ex:b2:structure-determination:ketone}
Uma substância desconhecida tem $M = 86$, uma banda de IV perto de \qty{1715}{cm^{-1}}, raias de carbono em
208.9, 45.7, 29.8, 17.4 e 13.7~ppm, e sinais de prótons em 2.41 (t, 2H), 2.13 (s, 3H), 1.61 (sexteto, 2H)
e 0.91 (t, 3H). Um grau de insaturação, uma cetona saturada (nenhum próton de aldeído perto de 9.7). Cinco
carbonos, cinco raias: nenhuma simetria. O singleto de três prótons em 2.13 é um \ce{CH3} vizinho da
carbonila; a sequência tripleto, sexteto, tripleto é um grupo propila cujo \ce{CH2} em 2.41 toca a
carbonila. Pentan-2-ona, \ce{CH3COCH2CH2CH3}. Seu isômero simétrico, a pentan-3-ona, mostraria três raias
de carbono e nenhum singleto; seu \omterm{def:b2:mass-spec-atomic:spectrum}{espectro de massas} tem o íon de McLafferty em 58 que falta à
pentan-3-ona (\cref{ch:b2:mass-spec-atomic}).
% ledger: c13:pentan-2-one, nmr:pentan-2-one, ir:CO-ketone, ms:pentan-2-one
\end{example}

\begin{example}[Um éster aromático, C$_9$H$_{10}$O$_2$]\label{ex:b2:structure-determination:ester}
Uma substância desconhecida de fórmula \ce{C9H10O2} (cinco graus de insaturação) mostra uma banda de
carbonila e sete raias de carbono: 166.5 (C), 132.8 (CH), 130.6 (C), 129.6 (CH), 128.3 (CH), 60.9
(\ce{CH2}) e 14.3 (\ce{CH3}). Seu espectro de prótons: 8.05 (m, 2H), 7.35--7.63 (m, 3H), 4.37 (q, 2H), 1.38
(t, 3H). Um anel benzênico monossubstituído (quatro raias de carbono \omterm{def:b2:huckel:aromatic}{aromático}, duas delas para dois
carbonos cada; cinco prótons \omterm{def:b2:huckel:aromatic}{aromáticos}) responde por quatro insaturações, a carbonila pela quinta. O par
quarteto--tripleto é um grupo etila; seu \ce{CH2} em 4.37 (e em 60.9 em carbono) está ligado a um
oxigênio. A carbonila de éster em 166.5 está ligada ao anel: os dois prótons em 8.05, vizinhos dela, são
deslocados para campo baixo. Benzoato de etila, \ce{C6H5COOCH2CH3}.
% ledger: c13:ethylbenzoate, nmr:ethylbenzoate
\end{example}

\section{As armadilhas}

\begin{itemize}
  \item \emph{Prótons trocáveis} (\ce{O-H}, \ce{N-H}) dão sinais largos em deslocamentos variáveis,
        muitas vezes não estão acoplados a seus vizinhos, e desaparecem quando se agita uma gota de
        \ce{D2O} com a amostra.
  \item \emph{Sinais sobrepostos}: dois carbonos de um anel benzênico, ou vários \ce{CH2} de uma cadeia,
        podem compartilhar uma raia; conte as raias em relação à fórmula antes de concluir sobre a simetria.
  \item \emph{Prótons diastereotópicos}: os dois prótons de um \ce{CH2} vizinho de um estereocentro não
        são equivalentes, podem ter deslocamentos diferentes e se acoplar entre si.
  \item \emph{Padrões de segunda ordem}: quando dois prótons acoplados têm deslocamentos mais próximos que
        algumas vezes sua constante de acoplamento, os multipletos se inclinam um para o outro e a regra
        do $n+1$ falha; um espectrômetro de campo mais alto restabelece padrões simples.
\end{itemize}

\begin{safety}{\ghs[5mm]{GHS02}\,\ghs[5mm]{GHS07}\,\ghs[5mm]{GHS08}}
A butan-2-ona é inflamável. O ímã de um espectrômetro de RMN está sempre ligado: nada de ferramentas de
aço, cilindros de gás ou cartões magnéticos perto dele, e nenhum acesso para pessoas com marca-passo.
% ledger: ghs:butanone, ghs:benzylacetate
\end{safety}

\section{Exercícios}

\begin{exercise}[$\star$]\label{exo:b2:structure-determination:1}
Quantas raias mostra o espectro de \ce{^{13}C} desacoplado de cada xileno (1,2-, 1,3- e
1,4-dimetilbenzeno)?
\end{exercise}

\begin{exercise}[$\star$]\label{exo:b2:structure-determination:2}
Preveja os espectros DEPT-135 e DEPT-90 do butan-2-ol.
\end{exercise}

\begin{exercise}[$\star$]\label{exo:b2:structure-determination:3}
A que tipo de carbono pertence mais provavelmente uma raia em 209, 170, 129, 64 e 14~ppm?
\end{exercise}

\begin{exercise}[$\star$]\label{exo:b2:structure-determination:4}
Calcule o grau de insaturação de \ce{C8H8O2} e proponha uma estrutura que contenha um anel benzênico.
\end{exercise}

\begin{exercise}[$\star\star$]\label{exo:b2:structure-determination:5}
Um líquido \ce{C4H8O} mostra uma banda de IV em \qty{1715}{cm^{-1}} e raias de carbono em 209, 37, 29 e 8~ppm
(DEPT-135: 37 para baixo, 29 e 8 para cima, 209 ausente). Identifique-o.
\end{exercise}

\begin{exercise}[$\star\star$]\label{exo:b2:structure-determination:6}
Como você distinguiria a pentan-2-ona da pentan-3-ona apenas por RMN de \ce{^{13}C}? Por RMN de prótons?
Por \omterm{def:b2:mass-spec-atomic:spectrum}{espectrometria de massas}?
\end{exercise}

\begin{exercise}[$\star\star$]\label{exo:b2:structure-determination:7}
Um éster tem $M = 88$ e sinais de prótons em 4.12 (q, 2H), 2.04 (s, 3H) e 1.26 (t, 3H) (dados do
exercício). Identifique-o.
\end{exercise}

\begin{exercise}[$\star\star$]\label{exo:b2:structure-determination:8}
Distinga o 1,2-diclorobenzeno do 1,4-diclorobenzeno por RMN de \ce{^{13}C}.
\end{exercise}

\begin{exercise}[$\star\star$]\label{exo:b2:structure-determination:9}
No 2-clorobutano, por que os dois prótons do grupo \ce{CH2} não são equivalentes? O que isso causa no
espectro de prótons?
\end{exercise}

\begin{exercise}[$\star\star\star$]\label{exo:b2:structure-determination:10}
Uma substância desconhecida \ce{C9H10O} mostra uma banda de IV em \qty{1690}{cm^{-1}}, raias de carbono em
200.8 (C), 137.0 (C), 132.9 (CH), 128.6 (CH), 128.0 (CH), 31.8 (\ce{CH2}) e 8.3 (\ce{CH3}), e sinais de
prótons em 7.95 (m, 2H), 7.40--7.55 (m, 3H), 2.98 (q, 2H) e 1.22 (t, 3H) (dados do exercício).
Identifique-a, e explique o baixo número de onda da carbonila.
\end{exercise}

\begin{exercise}[$\star\star\star$]\label{exo:b2:structure-determination:11}
Quatro isômeros \ce{C5H10O2}: ácido pentanoico, butanoato de metila, propanoato de etila e etanoato de
propila. Dê para cada um uma característica de seu espectro de IV ou de RMN de prótons que o distingue
dos outros três.
\end{exercise}

\begin{exercise}[$\star\star\star$]\label{exo:b2:structure-determination:12}
Um relatório atribui a raia em 166.5~ppm do benzoato de etila ao carbono do anel que leva o éster e a
raia em 130.6~ppm ao carbono carbonílico. Mostre, com os dados de \omterm{def:b2:structure-determination:dept}{DEPT} e o gráfico dos deslocamentos,
que a atribuição não pode estar certa.
\end{exercise}

\section{Problema: Quatro espectros do jasmim}

\begin{problem}[{Problema de fim de semana --- a fórmula a partir do íon molecular, a banda de carbonila
no infravermelho, o carbono 13 e o DEPT, o espectro de prótons, e uma estrutura confrontada com os
fragmentos}]
\label{pb:b2:structure-determination:1}
Um líquido incolor com odor floral de jasmim dá os seguintes dados. \omterm{def:b2:mass-spec-atomic:spectrum}{Espectro de massas} (NIST
WebBook): $m/z$ 150 (31.7), 151 (3.1), 108 (100), 91 (71.7), 90 (40.6), 43 (37.6). Infravermelho (dados do
exercício): bandas fortes em 1740 e \qty{1230}{cm^{-1}}, bandas fracas logo acima de \qty{3000}{cm^{-1}},
nenhuma banda larga perto de \qty{3300}{cm^{-1}}. Carbono 13, CDCl$_3$ (dados do PubChem): 170.70, 136.14,
128.56, 128.24, 66.24, 20.82~ppm. Prótons, CDCl$_3$, \qty{90}{MHz}: 7.33 (m, 5H), 5.09 (s, 2H), 2.06 (s, 3H).
% ledger: use:benzylacetate, ms:benzylacetate, c13:benzylacetate, nmr:benzylacetate, ir:CO-ester, iso:C12, iso:C13, mass:C12, mass:H1, mass:O16

\begin{center}
\begin{tikzpicture}
\pgfplotsset{dept/.style={width=0.8\linewidth, height=3.2cm, ycomb, x dir=reverse, ymin=0, ymax=1.15,
  xmin=0, xmax=200, ytick=\empty, axis y line=none, axis x line=bottom, x axis line style={-}, tick label style={font=\tiny},
  title style={font=\scriptsize, at={(0.02,0.82)}, anchor=west}, every axis plot/.append style={thick}}}
\begin{axis}[dept, title={desacoplado}, xtick={0,20,40,60,80,100,120,140,160,180,200}]
\addplot[omDef] table[x=shift, y=dec] {figdata/out/bachelor-2-nmr-spectra-c.dat};
\end{axis}
\end{tikzpicture}
\end{center}
% figdata: bachelor-2/nmr-spectra.py part c

\medskip
\noindent\textbf{Parte I --- A fórmula.}
\begin{enumerate}
  \item A partir da razão 151/150, estime o número de carbonos.
  \item O que diz a \omterm{def:b2:mass-spec-atomic:nitrogen-rule}{regra do nitrogênio}?
  \item Com nove carbonos, o que resta da massa 150? Proponha uma fórmula com oxigênio.
  \item Por que \ce{C10H14O}, também de massa nominal 150, é menos provável?
  \item Calcule o grau de insaturação.
  \item Calcule a \omterm{def:b2:mass-spec-atomic:exact-mass}{massa monoisotópica} da fórmula escolhida.
\end{enumerate}

\medskip
\noindent\textbf{Parte II --- O infravermelho.}
\begin{enumerate}[resume]
  \item Atribua a banda em \qty{1740}{cm^{-1}}. Que família sua posição sugere?
  \item Atribua a banda em \qty{1230}{cm^{-1}}.
  \item O que indicam as bandas fracas acima de \qty{3000}{cm^{-1}}?
  \item O que exclui a ausência de banda larga perto de \qty{3300}{cm^{-1}}?
  \item A carbonila é conjugada com o anel? Explique a partir de seu número de onda.
\end{enumerate}

\medskip
\noindent\textbf{Parte III --- Carbono 13 e \omterm{def:b2:structure-determination:dept}{DEPT}.}
\begin{enumerate}[resume]
  \item Atribua cada raia a um tipo de carbono.
  \item Que raia aponta para baixo em DEPT-135?
  \item Que raias desaparecem em DEPT-135?
  \item Que raias permanecem em DEPT-90?
  \item Quantos tipos de carbono tem um anel benzênico monossubstituído? Explique por que só aparecem
        duas raias CH \omterm{def:b2:huckel:aromatic}{aromáticas}.
  \item A raia em 128.24 é a mais alta. Ela representa o maior número de carbonos?
\end{enumerate}

\medskip
\noindent\textbf{Parte IV --- RMN de prótons e estrutura.}
\begin{enumerate}[resume]
  \item Atribua os três sinais de prótons.
  \item Por que os sinais em 5.09 e 2.06 são singletos?
  \item Por que o \ce{CH2} em 5.09 está tão deslocado para campo baixo?
  \item Monte a estrutura e dê seu nome.
  \item Seu isômero, o 2-feniletanoato de metila, \ce{C6H5CH2COOCH3}, tem a mesma fórmula. Que dado o
        exclui?
  \item Explique os fragmentos em 108, 91 e 43.
  \item Quantas raias mostraria um espectro de \ce{^{13}C} desacoplado perfeitamente resolvido do
        composto, e quantas delas apontam para baixo em DEPT-135?
\end{enumerate}
\end{problem}
